\section{Preliminares}
Nesta seção é descrito o ambiente de desenvolvimento, informando qual o hardware e softwares utilizados na implementação dos algoritmos e execução dos testes. Além disso, é apresentado a forma como o projeto foi organizado em diretórios e aquivos. Por fim, é feita algumas definições necessárias para um melhor entendimento das seções que se seguem.

\subsection{Ambiente de desenvolvimento}
Para a implementação dos algoritmos e execução dos testes foi utilizado um computador de mesa comum com o sistema operacional Ubuntu. Todos os algoritmos e funções auxiliares foram implementados com a linguagem de programação C. As linguagens de programação Bash e Python foram utilizadas de forma secundária para automatização dos testes e geração de gráficos, respectivamente. O hardware, softwares e versões utilizados são especificados na tabela \ref{tab:hardware-and-softwares}.

\begin{table}[H]
    \centering
    \Caption{\label{tab:hardware-and-softwares}Hardware e softwares utilizados}
    \begin{tabular}{ c | l | l }
        \hline
        Tipo                      & Componente                & Produto\\
        \hline
        \multirow{3}{*}{Hardware} & Placa mãe                 & MSI B450M PRO-VDH MAX\\
                                  & Processador               & AMD\textregistered \text{ }Ryzen 5 4600g\\
                                  & Memória principal         & 2x XPG 8GB DDR4\\
        \hline
        \multirow{4}{*}{Software} & Sistema operacional       & Linux Ubuntu 22.04 LTS\\
                                  & Editor de texto           & Microsoft Visual Studio Code\\
                                  & Compilador C              & GNU Compiler Collection GCC\\
                                  & Linguagens de programação & C, Bash e Python\\
        \hline
    \end{tabular}
\end{table}

\subsection{Organização do código fonte}
A figura \ref{fig:projeto-ordenacao} ilustra a forma como código fonte do projeto foi organizado, enquanto que a tabela \ref{tab:projeto-ordenacao} dá uma breve descrição de cada arquivo.

\begin{figure}[H]
\Caption{\label{fig:projeto-ordenacao}Árvore de diretórios e arquivos do projeto}
\centering
\begin{forest}
    for tree={font=\ttfamily, grow'=0,
    folder indent=.9em, folder icons,
    edge=densely dotted}
    [
        [graficos
            [graficos.ipynb, is file]
        ]
        [scripts
          [testes.sh, is file]
        ]
        [src
          [main.c, is file]
          [ordenacao.c, is file]
          [ordenacao.h, is file]
          [utils.c, is file]
          [utils.h, is file]
        ]
        [.gitignore, is file]
        [README.md, is file]
    ]
\end{forest}
\vspace{3pt}
\Fonte{Elaborado pelo autor. Acesso público em https://github.com/jbrenorv/ordenacao}
\end{figure}

\begin{table}[H]
    \centering
    \Caption{\label{tab:projeto-ordenacao}Descrição dos arquivos do projeto}
    \begin{tabular}{ c | l }
        \hline
        Arquivo        & Descrição\\
        \hline
        graficos.ipynb & Arquivo exportado do notebook criado no Google Colab\\
        \hline
        testes.sh      & Implementa script em Bash para execução automática dos testes\\
        \hline
        main.c         & Implementa a função principal\\
        \hline
        ordenacao.c    & Implementa os algoritmos de ordenação\\
        \hline
        ordenacao.h    & Declara os protótipos dos algoritmos de ordenação\\
        \hline
        utils.c        & Implementa as funções auxiliares\\
        \hline
        utils.h        & Declara protótipos das funções auxiliares\\
        \hline
        .gitignore     & Declara diretórios e arquivos que o Git deve ignorar\\
        \hline
        README.md      & Documenta o projeto com instruções de execução dos testes\\
        \hline
    \end{tabular}
\end{table}

\subsection{Tipos de vetores}
Além do tamanho do vetor, a forma como os elementos estão dispostos inicialmente pode interferir no tempo e na quantidade de operações que cada algoritmo executa. Pensando nisso, com o objetivo de analisar como cada algoritmo se comporta para diferentes tipos de entrada, foram considerados três tipos de vetores nos testes, como mostra a tabela \ref{tab:tipos-de-vetores}.

\begin{table}[H]
    \centering
    \Caption{\label{tab:tipos-de-vetores}Tipos de vetores}
    \begin{tabular}{ c | l }
        \hline
        Tipo   & Descrição\\
        \hline
        1      & Vetores já ordenados em ordem crescente\\
        \hline
        2      & Vetores ordenados em ordem decrescente\\
        \hline
        3      & Vetores gerados de forma pseudo-aleatória\\
        \hline
    \end{tabular}
\end{table}